% ===================================================================
\section{Related Literature}\label{sec:lit}
% ===================================================================
We organise the review around problems rather than individual papers, because our
contribution sits at the intersection of three otherwise separate strands: how causal
counterfactuals are built for a single aggregate unit, how the oil-price effect of
geopolitical shocks is conventionally measured, and how inference is conducted with one
treated unit and a short post-period. Tools borrowed from machine learning and time
series are introduced where they are used (Section~\ref{sec:methods} and the
appendices). Here we position the causal design.

\subsection{Building counterfactuals for a single aggregate unit}
The problem of constructing a counterfactual for one treated unit is what synthetic
control was invented to solve. The method originates with \citet{abadiegardeazabal2003}
and \citet{abadie2010}, who build the counterfactual as a convex combination of
untreated donors matched on pre-treatment outcomes and formalise permutation (placebo)
inference for the single-unit case. \citet{abadie2015} extend it to comparative
politics, and \citet{abadie2021} codifies the practical requirements we adopt: a long,
shock-free pre-window, donors clean of the treatment (the SCM analogue of SUTVA), and a
good pre-treatment fit as a precondition for credibility. A subsequent literature
relaxes the convex-hull restriction. The augmented SCM of \citet{benmichael2021} adds a
ridge bias-correction permitting limited extrapolation when the treated unit lies
outside the donor hull, \citet{doudchenko2016} embed SCM, difference-in-differences, and
regression in a common regularised-regression family that allows unconstrained weights,
and \citet{athey2021mc} recast the problem as matrix completion. A caution runs through
all of it. Convex SCM weights are \emph{non-unique} when donors are correlated
\citep{abadielhour2021}, so many weight vectors yield near-identical fit and individual
weights should not be over-interpreted. We use this entire toolbox by running an
ensemble that spans the convex, augmented, and regularised-regression cases, and we read
cross-model weight dispersion through the non-uniqueness result rather than as
instability.

\subsection{Where SCM has been applied, and the gap we occupy}
Despite this methodological maturity, SCM has rarely been applied to a daily market
price. Most peer-reviewed applications estimate effects on \emph{non-price} aggregates
such as GDP, trade, and employment, including under discrete geopolitical shocks.
\citet{gharehgozli2017} estimates the GDP cost of sanctions on Iran and
\citet{mancini2024} the trade effect of the 2022 Russia--Ukraine war. 

The closest
mechanical precedent to a price application is \citet{becker2021}, who apply SCM to
Austrian retail fuel prices with other European countries as donors. However, that design differs
from ours in three ways that matter, as its treatment is an endogenous regulation rather
than an exogenous shock, its outcome an administered retail price rather than a traded
benchmark, and its donors cross-country rather than cross-asset. Also, \citet{juarez2024} use
an oil-price shock as the treatment but on regional GDP, not on a price. 

The
intersection we occupy, which is SCM on a crude-benchmark \emph{price} under an exogenous
\emph{geopolitical} disruption with cross-asset donors, has no direct precedent, and
that is the gap this thesis addresses. In terms of existing analysis, because the Hormuz event is recent, the only
existing approach until now is institutional \citep{kiel2026,oies2026,worldbank2026,iea2026,
dallasfed2026}; they share the same structure, assuming a barrels-per-day shortfall,
propagate it through a structural model, and output a price path, and none applies
synthetic control to the observed price.

\subsection{Oil prices and geopolitical risk: the incumbent methods}
The methods conventionally used for oil-price effects answer narrower or different
questions than ours. The dominant approach is the structural VAR with identified supply,
demand, and inventory shocks \citep{kilian2009,baumeisterhamilton2019}, with
\citet{hamilton2009} the narrative anchor for 2007--08. These decompose price changes
into structural components but yield no counterfactual path and require identification
restrictions for which a chokepoint event offers no standard precedent. 

A natural
alternative method would regress oil prices on the Geopolitical Risk index of
\citet{caldara2022}, but this recovers the \emph{average} response to a unit of
geopolitical risk pooled across many events, a different object from the effect of one
specific disruption, and the index spikes endogenously to the very event we study, so it
can serve neither as our estimand nor as a clean donor. Short-window event studies
\citep{brownwarner1985} complete the incumbent set and give a magnitude but not the
multi-month path policy requires.

\subsection{Inference with one treated unit and a short post-period}
Our setting makes inference non-standard, and the recent literature counsels restraint
in two places. The first is inference itself: permutation and placebo tests
\citep{abadie2010,abadie2015} are the field standard, with validity resting on an
approximate-symmetry assumption \citep{canay2017}. \citet{hahnshi2017} show that for
convex SCM this effectively requires Gaussian errors, which bounds how far placebo
$p$-values can be pushed for our nonlinear ensemble members, and \citet{chenyan2023}
formalise the in-time placebo into a mixed test with a permutation $p$-value. 

The second
is pretesting. \citet{roth2022} shows that pre-trend tests have low power at applied
sample sizes and can manufacture false confidence, and \citet{rambachanroth2023} propose
bounding rather than testing such violations, so we accordingly adopt a
report-and-interpret stance toward pre-period diagnostics rather than imposing pass/fail
thresholds. 

In our case, the value of a donor depends on how reliably it tracks the common
factor. A literature on matching reliability before differencing shows that a good
pre-period fit can \emph{mask} rather than remove contamination, and this is the logic we use to
split donors by pre-window factor $R^2$ when signing the contamination bias
(Section~\ref{sec:results}). Multiple testing across donors is controlled by
Benjamini--Hochberg \citep{benjaminihochberg1995}.

